\documentclass[10pt, a4paper]{article}
\usepackage{preamble}

\title{Mathematical Methods II \\
    \large Problem Class}
\author{Luke Phillips}
\date{October 2025}

\begin{document}

\maketitle

\newpage

\tableofcontents

\newpage

\section{Problems Class \texorpdfstring{$1$}{}}

\begin{problem}[$2012$ question $2$]
    Calculate $\int_{C}\vec{f}\cdot\hvec{t}\,d\ell$ where $\vec{f} = -x\vec{e}_1 + 4y\vec{e}_2$ and $C$ is the part of $x ^ 2 + 4y ^ 2 = 4$ from $(2, 0)$ to $(0, 1)$.

    \begin{solution}
        Parameterisation:
        \[
        \vec{x}(t) = 2\cos(t)\vec{e}_1 + \sin(t)\vec{e}_2\qquad\text{for } t \in [0, \frac{\pi}{2}].
        \]
        \begin{align*}
            \int_{C}\vec{f}\cdot\hvec{t}\,d\ell &= \int_{0}^{\frac{\pi}{2}}\vec{f}\cdot\hvec{t}\left|\frac{d\vec{x}}{dt}\right|\,dt \\
            &= \int_{0}^{\frac{\pi}{2}}\vec{f}\cdot\frac{d\vec{x}}{dt}\,dt \\
            &= \int_{0}^{\frac{\pi}{2}}\vec{f}\cdot(-2\sin(t)\vec{e}_1 + \cos(t)\vec{e}_2)\,dt \\
            &= \int_{0}^{\frac{\pi}{2}}(-x\vec{e}_1 + 4y\vec{e}_2)\cdot(-2\sin(t)\vec{e}_1 + \cos(t)\vec{e}_2)\,dt \\
            &= \int_{0}^{\frac{\pi}{2}}(-2\cos(t)\vec{e}_1 + 4\sin(t)\vec{e}_2)\cdot(-2\sin(t)\vec{e}_1 + \cos(t)\vec{e}_2)\,dt \\
            &= \int_{0}^{\frac{\pi}{2}}(4\cos(t)\sin(t) + 4\sin(t)\cos(t))\,dt \\
            &= 8\int_{0}^{\frac{\pi}{2}}\cos(t)\sin(t)\,dt \\
            &= 8\int_{1}^{0}(-u)\,du \\
            &= -8\left[\frac{u ^ 2}{2}\right]_{1}^{0} = 4.
        \end{align*}
    \end{solution}
\end{problem}


\begin{problem}[$2015$ question $8$]
    Calculate the line integral of $\vec{f}(\vec{x}) = z ^ 2\vec{e}_1 + x ^ 2\vec{e}_2 + y ^ 2\vec{e}_3$ along each of the following curves.

    \begin{enumerate}[label = (\alph*)]
        \item The straight line $C_1$ from $(0, 0, 0)$ to $(1, 1, 1)$.

        \item The circle $C_2$ with radius $a > 0$ in the $z = 0$ plane centred at $(1, 1, 0)$,
        traversed anti-clockwise.

        \item The triangle $C_3$ with vertices $(1, 1, 0) \to (1 + b, 1, 0) \to (1, 1 + b, 0)$ for $b > 0$.
    \end{enumerate}

    \begin{solution}
        \begin{enumerate}[label = (\alph*)]
            \item Need parameterisation:
            \[
            \vec{x}(t) = t\vec{e}_1 + t\vec{e}_2 + t\vec{e}_3\quad\text{for } t \in [0, 1]
            \]
            \[
            \frac{d\vec{x}}{dt} = \vec{e}_1 + \vec{e}_2 + \vec{e}_3.
            \]

            So
            \begin{align*}
                \int_{C_1}\vec{f}\cdot\hvec{t}\,d\ell &= \int_{0}^{1}(z ^ 2\vec{e}_1 + x ^ 2\vec{e}_2 + y ^ 2\vec{e}_3)\cdot(\vec{e}_1 + \vec{e}_2 + \vec{e}_3)\,dt \\
                &= \int_{0}^{1}(t ^ 2\vec{e}_1 + t ^ 2\vec{e}_2 + t ^ 2\vec{e}_3) \cdot (\vec{e}_1 + \vec{e}_2 + \vec{e}_3)\,dt \\
                &= \int_{0}^{1}3t ^ 2\,dt = \left[t ^ 3\right]_{0}^{1} = 1.
            \end{align*}

            \item Parameterisation for $C_2$:
            \[
            \vec{x}(t) = (1 + a\cos(t))\vec{e}_1 + (1 + a\sin(t))\vec{e}_2\quad\text{for } [0, 2\pi].
            \]
            \[
            \frac{d\vec{x}}{dt} = -a\sin(t)\vec{e}_1 + a\cos(t)\vec{e}_2.
            \]
            So
            \begin{align*}
                \int_{C_2}\vec{f}\hvec{t}\,d\ell &= \int_{0}^{2\pi}(z ^ 2\vec{e}_1 + x ^ 2\vec{e}_2 + y ^ 2\vec{e}_3) \cdot (-a\sin(t)\vec{e}_1 + a\cos(t)\vec{e}_2)\,dt \\
                &= \int_{0}^{2\pi}\left[(1 + a\cos(t)) ^ 2\vec{e}_2 + (1 + a\sin(t)) ^ 2\vec{e}_3\right]\cdot\left(-a\sin(t)\vec{e}_1 + a\cos(t)\vec{e}_2\right)\,dt \\
                &= \int_{0}^{2\pi}a\cos(t)(1 + a\cos(t)) ^ 2\,dt \\
                &= \int_{0}^{2\pi}a\cos(t)(1 + 2a\cos(t) + a ^ 2\cos ^ 2(t))\,dt \\
                &= a\int_{0}^{2\pi}\cos(t)\,dt + 2a ^ 2\int_{0}^{2\pi}\cos ^ 2(t)\,dt + a ^ 3\int_{0}^{2\pi}\cos ^ 3(t)\,dt \\
                &= 0 + 2\pi a ^ 2 + 0 = 2\pi a ^ 2.
            \end{align*}

            \item
            Split into three straight lines that are differentiable:
            \[
            \vec{x}(t) = \begin{cases}
                (1 + bt)\vec{e}_1 + \vec{e}_2 &\text{for } t \in [0, 1] \\
                (1 + b(1 -t))\vec{e}_1 + (1 + bt)\vec{e}_2 &\text{for } t \in [0, 1] \\
                \vec{e}_1 + (1 + b(1 -t))\vec{e}_2 &\text{for } t \in [0, 1].
            \end{cases}
            \]
            So
            \begin{align*}
                \oint_{C}\vec{f}\cdot\hvec{t}\,d\ell &= \int\vec{f}\cdot\hvec{t}\,d\ell + \int\vec{f}\cdot\hvec{t}\,d\ell + \int\vec{f}\cdot\hvec{t}\,d\ell \\
                &= \int_{0}^{1}(x ^ 2\vec{e}_2 + y ^ 2\vec{e}_3)\cdot(b\vec{e}_1)\,dt + \int_{0}^{1}(x ^ 2\vec{e}_2 + y ^ 2\vec{e}_3)\cdot(-b\vec{e}_1 + b\vec{e}_2)\,dt + \int_{0}^{1}(x ^ 2\vec{e}_2 + y ^ 2\vec{e}_3)\cdot(-b\vec{e}_2)\,dt \\
                &= \int_{0}^{1}0\,dt + \int_{0}^{1}bx ^ 2\,dt - \int_{0}^{1}bx ^ 2\,dt \\
                &= \int_{0}^{1}b(1 + b(1 - t)) ^ 2\,dt - \int_{0}^{1}b\,dt \\
                &= \int_{0}^{1}b(1 + b ^ 2 + b ^ 2t ^ 2 + 2b - 2bt)\,dt - \int_{0}^{1}b\,dt \\
                &= \left[b(1 + b ^ 2)t + \frac{b ^ 3}{3}t ^ 3 + 2b ^ 2t - b ^ 2t ^ 2 - b ^ 3t ^ 2\right]_{0}^{1} - \left[bt\right]_{0}^{1} \\
                &= b(1 + b ^ 2) + \frac{b ^ 3}{3} + 2b ^ 2 - b ^ 2 - b ^ 3 - b \\
                &= b ^ 2\left(\frac{b}{3} + 1\right).
            \end{align*}
        \end{enumerate}
    \end{solution}
\end{problem}

\begin{problem}[August $2024$ question $6$]
    The position vector in
    "ellipsoidal coordinates"
    is given by
    \[
    \vec{x}(r, \theta, \phi) = ar\sin(\theta)\cos(\phi)\vec{e}_1 + br\sin(\theta)\sin(\phi)\vec{e}_2 + cr\cos(\theta)\vec{e}_3
    \]
    where $r > 0$,
    $\theta \in [0, \pi]$,
    $\phi \in [0, 2\pi]$ and $a, b, c \in \R$ are given positive constants.

    Determine whether or not this $(r, \theta, \phi)$ coordinate system is orthogonal.

    \begin{solution}
        Calculate tangent vectors:
        \begin{align*}
            \pd[\vec{x}]{r} &= a\sin(\theta)\cos(\phi)\vec{e}_1 + b\sin(\theta)\sin(\phi)\vec{e}_2 + c\cos(\theta)\vec{e}_3 \\
            \pd[\vec{x}]{\theta} &= ar\cos(\theta)\cos(\phi)\vec{e}_1 + br\cos(\theta)\sin(\phi)\vec{e}_2 - cr\sin(\theta)\vec{e}_3
        \end{align*}
        then dotting them together
        \[
        \pd[\vec{x}]{r} \cdot \pd[\vec{x}]{\theta} \neq 0.
        \]
    \end{solution}
\end{problem}

\newpage

\section{Problems Class \texorpdfstring{$2$}{}}

\begin{problem}[$2013$ question $5$]
    Compute the surface integral $\int_{S}\vec{f}\cdot\hvec{n}\,dS$ of vector field $\vec{f} = (y, -x, x)$ over the surface $z = xy - \cos{x}$ for $x \in [0, 1]$,
    $y \in [0, 1]$,
    with upward normal.

    \begin{solution}
        \[
        \vec{f}(\vec{x}) = y\vec{e}_1 + x\vec{e}_2 + x\vec{e}_3.
        \]

        \begin{align*}
            \int_{S}\vec{f}\cdot\hvec{n}\,dS &= \int_{U}\vec{f}(\vec{x}(u, v)\cdot\hvec{n}\left|\pd[\vec{x}]{u}\times \pd[\vec{x}]{v}\right|\,dudv \\
            &= \int_{U}\vec{f}(\vec{x}(u, v)\cdot\left(\pd[\vec{x}]{u}\times \pd[\vec{x}]{v}\right)\,dudv.
        \end{align*}
        Parametrise the surface:
        \[
        z = xy - \cos{x}
        \]
        \[
        \vec{x}(x, y) = x\vec{e}_1 + y\vec{e}_2 + (xy - \cos{x})\vec{e}_3
        \]
        $x \in [0, 1]$,
        $y \in [0, 1]$.
        So
        \begin{align*}
            \pd[\vec{x}]{x} &= \vec{e}_1 + (y + \sin{x})\vec{e}_3 \\
            \pd[\vec{x}]{y} &= \vec{e}_2 + x\vec{e}_3
        \end{align*}
        \[
        \pd[\vec{x}]{x} \times \pd[\vec{x}]{y} = -(y + \sin{x})\vec{e}_1 + -x\vec{e}_2 + \vec{e}_3.
        \]
        Then
        \begin{align*}
            \int_{S}\vec{f}\cdot\hvec{n}\,dS &= \int_{0}^{1}\int_{0}^{1}(y\vec{e}_1 - x\vec{e}_2 + x\vec{e}_3)\cdot(-(y + \sin{x})\vec{e}_1 - x\vec{e}_2 + \vec{e}_3)\,dxdy \\
            &= \int_{0}^{1}\int_{0}^{1}(-y ^ 2 -y\sin{x} + x ^ 2 + x)\,dxdy \\
            &= \int_{0}^{1}\left[-y ^ 2x + y\cos{x} + \frac{1}{3}x ^ 3 + \frac{1}{2}x ^ 2\right]_{0}^{1}\,dy \\
            &= \int_{0}^{1}-y ^ 2 + y\cos(1) + \frac{5}{6} - y)\,dy \\
            &= \int_{0}^{1}-y ^ 2 + y(\cos(1) - 1) + \frac{5}{6})\,dy \\
            &= \left[-\frac{1}{3}y ^ 3 + \frac{1}{2}y ^ 2(\cos(1) - 1) + \frac{5}{6}y\right]_{0}^{1} \\
            &= -\frac{1}{3} + \frac{1}{2}\cos(1) + \frac{5}{6} - \frac{1}{2} = \frac{1}{2}\cos(1).
        \end{align*}
    \end{solution}
\end{problem}

\begin{problem}[$2020$ question $5$(a)]
    Calculate the surface integral $\oint_{S}\vec{f}\cdot\hvec{n}\,dS$ where $\vec{f}(\vec{x}) = \vec{x}$ and $S$ is the surface of the solid cylinder $\{x ^ 2 + y ^ 2 \leq 1, 0 \leq z \leq 4\}$,
    with outward normal.

    \begin{solution}
        \[
        \vec{f}(\vec{x}) = \vec{x} = x\vec{e}_1 + y\vec{e}_2 + z\vec{e}_3.
        \]

        $S$ is not smooth so divide into three pieces.
        $S_1$ -
        top,
        $S_2$ -
        bottom,
        $S_3$ -
        curved surface.

        Parametrise in cylindrical coordinates:
        \[
        \vec{x}(r, \theta, z) = r\cos(\theta)\vec{e}_1 + r\sin(\theta)\vec{e}_2 + z\vec{e}_3.
        \]
        $S_1$:
        \[
        \vec{x}(r, \theta) = r\cos(\theta)\vec{e}_1 + r\sin(\theta)\vec{e}_2 + 4\vec{e}_3\qquad\text{for } r \in [0, 1], \theta \in [0, 2\pi].
        \]
        \begin{align*}
            \int_{S_1}\vec{f}\cdot\hvec{n}\,dS &= \int_{S_1}\vec{f}\cdot\vec{e}_3\,dS \\
            &= \int_{0}^{2\pi}\int_{0}^{1}\vec{x}\cdot\vec{e}_3\left|\pd[\vec{x}]{r}\times\pd[\vec{x}]{\theta}\right|\,drd\theta \\
            &= \int_{0}^{2\pi}\int_{0}^{1}zr\,drd\theta \\
            &= 4\int_{0}^{2\pi}\int_{0}^{1}r\,drd\theta \\
            &= 4\pi.
        \end{align*}

        $S_2$:
        \[
        \vec{x}(r, \theta) = r\cos(\theta)\vec{e}_1 + r\sin(\theta)\vec{e}_2\qquad\text{for } r \in [0, 1], \theta \in [0, 2\pi].
        \]
        \begin{align*}
            \int_{S_2}\vec{f}\cdot\hvec{n}\,dS &= \int_{S_2}\vec{f}\cdot(-\vec{e}_3)\,dS \\
            &= \int_{0}^{2\pi}\int_{0}^{1}-zr\,drd\theta \\
            &= 0 &&(\text{z = 0}).
        \end{align*}


        $S_3$:
        \[
        \vec{x}(\theta, z) = \cos(\theta)\vec{e}_1 + \sin(\theta)\vec{e}_2 + z\vec{e}_3\quad\text{for } \theta \in [0, 2\pi], z \in [0, 4].
        \]
        \begin{align*}
            \int_{S_3}\vec{f}\cdot\hvec{n}\,dS &= \int_{S_3}\vec{f}\cdot\vec{e}_r\,dS
            \intertext{Either convert $\vec{e}_r$ to Cartesian or $\vec{f}$ to cylindrical}
            &= \int_{S_3}(r\vec{e}_r + z\vec{e}_z) \cdot \vec{e}_r\,dS \\
            &= \int_{0}^{4}\int_{0}^{2\pi}r\,d\theta dz \\
            &= 8\pi.
        \end{align*}
        So overall,
        \[
        \oint_{S}\vec{f}\cdot\hvec{n}\,dS = \int_{S_1}\vec{f}\cdot\hvec{n}\,dS + \int_{S_2}\vec{f}\cdot\hvec{n}\,dS + \int_{S_3}\vec{f}\cdot\hvec{n}\,dS = 12\pi
        \]
    \end{solution}
\end{problem}

\begin{problem}[August $2025$ question $2$(b)]
    Compute the volume of a solid chocolate egg given by
    \[
    \{r < a\cos ^ 3{\theta},\quad\theta \in [0, \pi / 2]\}
    \]
    in spherical coordinates,
    where $a > 0$ is constant.

    \begin{solution}
        Parametrise in spherical coordinates
        \[
        \vec{x}(r, \theta, \phi) = r\sin(\theta)\cos(\phi)\vec{e}_1 + r\sin(\theta)\sin(\phi)\vec{e}_2 + r\cos(\theta)\vec{e}_2
        \]

        \begin{align*}
            |V| &= \int_{V}\,dV \\
            &= \int_{0}^{2\pi}\int_{0}^{\pi / 2}\int_{0}^{a\cos ^ 3\theta}h_rh_{\theta}h_{\phi}\,drd\theta d\phi \\
            &= \int_{0}^{2\pi}\int_{0}^{\pi / 2}\int_{0}^{a\cos ^ 3\theta}r ^ 2\sin(\theta)\,drd\theta d\phi \\
            &= \int_{0}^{2\pi}\int_{0}^{\pi / 2}\left[\frac{1}{3}r ^ 3\sin(\theta)\right]_{0}^{a\cos ^ 3\theta}\,d\theta d\phi \\
            &= \frac{a ^ 3}{3}\int_{0}^{2\pi}\int_{0}^{\pi / 2}\cos ^ 9(\theta)\sin(\theta)\,d\theta d\phi \\
            &= \frac{a ^ 3}{3}\int_{0}^{2\pi}\left[-\frac{1}{10}\cos ^ {10}(\theta)\right]_{0}^{\pi / 2}\,d\phi \\
            &= \frac{a ^ 3}{30}\int_{0}^{2\pi}\,d\phi = \frac{\pi a ^ 3}{15}.
        \end{align*}
    \end{solution}
\end{problem}












\end{document}